% Edit this file for each poster. All text below is illustrative, not research evidence.
\newcommand{\LeftContent}{
{\PosterPaperTitleFont\bfseries Your paper title goes here\par}
\vspace{0.30in}
{\PosterAuthorFont Author One, Author Two\par}
{\PosterAuthorFont Mila / McGill University\par}
\PosterSection{Summary}{\PosterBullets{
\item Help a reader understand the work without talking to you.
\item Summarize the question, contribution, and main finding.
\item Give a \textbf{quick overview} worth saving in a photo.
}}
\PosterSection{Method}{\PosterBullets{
\item Explain the essential idea in a few short bullets.
\item Use a small diagram when it makes the method clearer.
\item Keep this pane readable in about one minute.
}}
\vspace{0.35in}
\FigurePlaceholder{5in}{Method diagram\\[0.15in]Input $\rightarrow$ Model $\rightarrow$ Output}
\PosterSection{Contact}{\href{mailto:you@example.org}{you@example.org}\\Conference name / Year}
}
\newcommand{\CenterContent}{
{\color{PosterPurple}\PosterTakeawayFont\bfseries Your main\\takeaway.\par}
\vspace{0.45in}
{\PosterSubtitleFont One or two sentences for people walking past.\par}
\vspace{0.65in}
\FigurePlaceholder{10in}{Optional main-result figure\\[0.25in]Use only if it makes the takeaway clear at a glance.}
\vspace{0.20in}
{\PosterBodyFont\textbf{Figure 1.} Give just enough context to understand the result at a glance.\par}
\PosterCenterSection{Why it matters}{State the significance in plain language. Keep the takeaway dominant.}
}
\newcommand{\RightContent}{
\PosterSection{Experiments and details}{\PosterBullets{
\item Choose figures and tables to support a conversation.
\item Use short captions; explain the details in person.
\item Keep labels, units, and evaluation conditions clear.
}}
\vspace{0.25in}
\FigurePlaceholder{5in}{Discussion figure or table}
\PosterSection{Limitations}{Note the assumption or limitation needed to interpret the results.}
\PosterSection{Next steps}{Add only details that help explain your main points.}
\PosterSection{References}{\PosterReferences{
[1] A. Author et al. Paper title. Venue, year.\\[0.15in]
[2] B. Author et al. Related work. Venue, year.
}}
\PosterSection{Acknowledgments}{\PosterReferences{Brief funding and collaborator credits. Use the logos of your actual affiliations.}}
}
% Generate locally: python3 scripts/make-qr.py URL my-qr.png --mode light
% Light: foreground #662E7D, background #FFFFFF.
% Dark: foreground #FFFFFF, background #662E7D. Scan-test both codes.
% Replace both QR images with codes for your actual paper/code URLs.
% Example destinations: https://example.org and https://example.org/code .
\newcommand{\PaperQRFile}{assets/qrcodes/example-qr.png}
\newcommand{\CodeQRFile}{assets/qrcodes/example-code-qr.png}
\newcommand{\MilaLogoFile}{assets/logos/mila-color.pdf}
\newcommand{\McgillLogoFile}{assets/logos/mcgill-color.pdf}
\newcommand{\IclrLogoFile}{assets/logos/iclr-color.png}
\newcommand{\FooterContent}{\PosterFooterLayout}
